\subsection{Global congruences on the zero free elliptic slice}

\begin{statement}{E.96}{Proved}{An explicit global class-two locus and its exact modulo-three inertia image}
Keep the cubic curve, rational base tangent, genuine arithmetic action,
and pro-three groups of E.37. Use the actual change of marking of E.95:
the full action is transported by $\Psi^{-1}\rho\Psi$, every quotient
map by $F\Psi$, and every path by the same isomorphism. Its Tate matrix
fixes the eight peripheral vectors and commutes with the displayed
first-degree arithmetic matrices. The following formulas use the
resulting phase-zero path constants.

Let $W\subset V$ be the rank-eight toric lattice and put
$M=\bigwedge^2V$, $N_2=N/\gamma_3N$. Define the following actual
subset of global continuous cohomology:
\[
 \mathcal H_0=\ker\left(
 H^1(\mathbb Q,V)\longrightarrow H^1(\mathbb Q,T_3E)
       =\mathbb Z_3t_1\oplus C_3t_2
       \longrightarrow\mathbb Z_3t_1\right).
 \tag{E96.1}
\]
The arrows are induced by the compact projection and by E.37's
specified coordinates. Every element of $\mathcal H_0$ has unique
coordinates $(\epsilon,\beta)$, with $\epsilon\in\{0,1,2\}$ and
$\beta\in H^1(\mathbb Q,W)$, represented by
\[
 \alpha=\epsilon u_{f_2}+\kappa(\beta).
 \tag{E96.2}
\]
Here $u_{f_2}(\tau)=f_2$, its value at $\tau'$ is $Jf_2$,
and its value at $c$ is zero. The toric group is the actual Kummer
group
$(\widehat{\mathbb Q^\times}^{\,3})^2
 \times(\widehat{K^\times}^{\,3})^3$ of E.92;
its addition with the representative $\epsilon$ retains the
finite carry $c_T$ in that statement.

Use the order $(f_1,f_2,A_0,A_1,A_2,B_0,B_1,B_2,C_1,C_2)$.
The following fixed columns occur in the formula:
\[
 \begin{aligned}
 h_*&=\tfrac12(I+J)f_1-q_0,\\
 q_0&=(0,0,0,0,0,0,\tfrac14,-\tfrac14,\tfrac14,-\tfrac14)^t,\\
 c_3&=(0,0,1,1,1,2,\tfrac12,\tfrac12,\tfrac32,\tfrac32)^t,\\
 c_\zeta&=(0,0,0,1,-1,0,\tfrac12,-\tfrac12,
                                  \tfrac12,-\tfrac12)^t.
 \end{aligned}
 \tag{E96.3}
\]
On $G_K$, the literal toric carry satisfies
$\kappa(c_T)=c_3\kappa_3+c_\zeta\kappa_\zeta$, where
$\kappa_\zeta(g)=(\chi(g)-1)/3$ and $\kappa_3$ is the
unchanged ordinary Kummer cocycle of three. E.95's eight actual
framed-point cocycles have the specified rational columns
\[
 k_d(g)=(A_g-I)b_d+r_d\kappa_3(g)
                          +s_d(\chi(g)-1)/3.
 \tag{E96.4}
\]
All $3b_d,3r_d,3s_d$ are integral over $\mathbb Z_3$,
and $r_d,s_d\in W\otimes\mathbb Q_3$. These are the
explicit columns retained with the E.95 path calculation.

The quotient lattice used here is exactly $M/\mathbb Z_3w$,
with $w$ from E.37. Let $p$ denote the coordinate
$e_{A_0}\wedge e_{A_2}$, whose coefficient in $w$ is one.
Define the integral quotient map and its integral section by
\[
 Qm=(m-m_pw)_{i\ne p}\in\mathbb Z_3^{44},\qquad
 Sx=\text{the vector obtained by inserting zero at }p.
 \tag{E96.5}
\]
Thus $QS=I$, $SQ=I-we_p^t$, and $Qw=0$.

For each prime $\mathfrak q\nmid3$ of $K$, use the compatible
Frobenius and tame generator of E.37 and put
\[
 \begin{gathered}
 q=N\mathfrak q,\qquad k=q-1,\qquad n=k/3,\qquad
 t=\kappa_3(F_{\mathfrak q}),\\
 y=\kappa(\beta)(F_{\mathfrak q})\in W,\qquad
 v=\alpha(\tau_{\mathfrak q})
     =\sum_dv_de_d\in W,\\
 L_b(v)=\sum_dv_db_d\wedge e_d,\qquad
 L_r(v)=\sum_dv_dr_d\wedge e_d,\qquad
 L_s(v)=\sum_dv_ds_d\wedge e_d.
 \end{gathered}
 \tag{E96.6}
\]
The valuations defining $v$ and the Kummer values defining $y,t$
use the actual field classes and source marking. Define the
following integral column of forty-four coordinates:
\[
 \begin{split}
 \mathcal N_{\mathfrak q}=Q\bigl(&
   \epsilon k\,h_*\wedge v-3kL_b(v)+3y\wedge v
        +\epsilon t\,c_3\wedge v\\
   &+\epsilon n\,c_\zeta\wedge v
        -3tL_r(v)-kL_s(v)\bigr).
 \end{split}
 \tag{E96.7}
\]
In particular the elliptic Frobenius matrix has disappeared
from this exact formula. The original E.37 correction is
\[
 c_{\mathfrak q}=-\frac{\mathcal N_{\mathfrak q}}{3(q-1)}.
 \tag{E96.8}
\]
Consequently the entire image of
$H^1(\mathbb Q,N_2)\to H^1(\mathbb Q,V)$ inside
$\mathcal H_0$ is the following explicitly specified set:
\[
 \begin{gathered}
 \mathcal N_{\mathfrak q}\in
       3^{v_3(q-1)+1}\mathbb Z_3^{44}\quad
                         \text{for every }\mathfrak q\nmid3,\\
 \text{for every }m\ge0,\quad
 \mathcal N_{\mathfrak q}\in
       3^{v_3(q-1)+1+m}\mathbb Z_3^{44}\quad
                         \text{outside a finite set of primes.}
 \end{gathered}
 \tag{E96.9}
\]
On the subgroup whose toric valuation vectors have a fixed finite
support, the second line is automatic and the first line has
only that finite support to check. Formula (E96.7), or its
forty-four expanded polynomial rows in the certificate,
evaluates every congruence.

There is also a small exact quotient of this image. Put
\[
 D_0=O=C_0,\qquad D_1=A_0,\qquad D_2=B_0,\qquad
 e_O=-\sum_{d\ne O}e_d.
 \tag{E96.10}
\]
On the stratum with compact coordinate $\epsilon t_2$,
the image on all prime inertia vectors modulo three is
exactly
\[
 \bigoplus_{p\ne3}\mathbb F_3e_{D_\epsilon}.
 \tag{E96.11}
\]
The sum is indexed by rational primes. At every prime of $K$
above $p$ use its corresponding valuation vector, with the
conjugate-prime equivariance induced by the actual torus.
In the displayed eight coordinates this says respectively
\[
 \overline v_{\mathfrak q}\in
 \begin{cases}
   \mathbb F_3(1,1,1,1,1,1,1,1),&\epsilon=0,\\
   \mathbb F_3(1,0,0,0,0,0,0,0),&\epsilon=1,\\
   \mathbb F_3(0,0,0,1,0,0,0,0),&\epsilon=2.
 \end{cases}
 \tag{E96.12}
\]
Actual rational tangential sections at the three indicated
cusps, twisted by rational Kummer classes, give all elements
of (E96.11). The pure toric class $\kappa(7)e_{A_0}$ is
excluded on the $\epsilon=0$ stratum; the two classes obtained
by twisting the actual $A_0$ and $B_0$ tangential sections by
$\kappa(7)$ have zero correction at the marked prime of norm seven.

Every nonempty full class-two fibre over a point of (E96.9)
is the torsor of E.37 under
\[
 \mathcal R_{\{p\ne3\}}(\mathbb Z_3)
       \oplus\mathbb Z_3^{24}\oplus C_3^9\oplus C_9^5.
 \tag{E96.13}
\]
E.93 gives its exact central localization map. The slice
with nonzero free elliptic coordinate, further geometric
quotients, and the quartic curve remain in the scope of
the active E.23 task.

\emph{Manuscript: section51.16, folios733v--736v,
scans172--178, formulas(371)--(379); Part E(c),(d)
and follow-up(4).}
\end{statement}
\begin{proof}
\emph{Coordinates and the actual marking.}
The compact projection and the vanishing of its global
degree-zero group give the injective toric Kummer map
used in E.92. Its description of the compact cohomology
then gives the unique representatives (E96.2). The
normalization $u_{f_2}(c)=0$ and the chosen two free
generators specify an actual cocycle. In particular this
choice uses the finite carry rather than treating the
three possible representatives as an additive direct
summand.

The E.95 change of marking has an actual lift $\Psi$
on the free pro-three group. Sending a cocycle $z(g)$
to $\Psi^{-1}z(g)$ is a bijection of continuous cocycles
for the actions $\rho$ and $\Psi^{-1}\rho\Psi$, and sends
geometric conjugations to geometric conjugations. Its
inverse is induced by $\Psi$. The lower central series
is characteristic, so this isomorphism descends to
$N_2$ and to its abelian quotient. On the latter it
fixes every $e_d$; on the compact Tate module it is the
scalar $u^{-1}$ with $u\equiv1\pmod3$. Thus it preserves
the zero free-coordinate slice and its compact residue
$\epsilon$. All obstruction and inertia assertions
therefore transfer through this specified isomorphism.
The actual transported peripheral loop at $O$ still
has abelian value $e_O$ in (E96.10).

E.92 and the full cochain comparison in E.95 give
\[
 3u_{f_2}(g)=c_3\kappa_3(g)
       +c_\zeta(\chi(g)-1)/3+(A_g-I)h_*.
 \tag{E96.14}
\]
This is an identity on the actual arithmetic source.
The two sides are cocycles, agree on both free
generators, and respect complex conjugation. The
columns in (E96.3) are those obtained by the literal
Shapiro map on $c_T$, including its third roots of
unity. Equation (E96.4) has the same full-source
meaning by the framed-point construction of E.95.

\emph{The exact inverse of the local differential.}
Fix $\mathfrak q\nmid3$. On $W$ the Frobenius matrix
is $qI$. For $d=q^{-1}\bigwedge^2A-I$, the following
identities hold over $\mathbb Q_3$:
\[
 d(h\wedge e_a)=(A-I)h\wedge e_a,
 \qquad d|_{\wedge^2W}=(q-1)I.
 \tag{E96.15}
\]
They follow by applying the exterior square to the
two factors; no splitting of the integral extension
$0\to W\to V\to T_3E\to0$ is used.
Write
\[
 h=\epsilon h_*/3,\qquad
 u=y+\epsilon(c_3t+c_\zeta k/3)/3.
 \tag{E96.16}
\]
Equation (E96.14) says exactly that
$\alpha(F)=(A-I)h+u$. Define the rational exterior column
\[
 X=h\wedge v-L_b(v)
          +\frac{u\wedge v-tL_r(v)-(k/3)L_s(v)}{k}.
 \tag{E96.17}
\]
Apply (E96.15) to every term. Substitution of (E96.4)
then gives the literal equality
\[
 dX=\alpha(F)\wedge v
          -\sum_av_a k_a(F)\wedge e_a
       =D_{\mathfrak q}.
 \tag{E96.18}
\]
The last column is the actual local extension defect
in E37.10, including its peripheral sign.

The coefficient $w_p=1$ proves all three identities
following (E96.5). It also proves that $\ker Q$ on
the integral lattice is exactly $\mathbb Z_3w$:
$Qm=0$ gives $m=m_pw$ with integral $m_p$.
Thus $Q$ is the full quotient map to the actual
domain of $d_{\mathfrak q}$. Its rational kernel is
$\mathbb Q_3w$. By E.37, the induced differential
from this quotient to $M_0$ is invertible over
$\mathbb Q_3$. Equation (E96.18) therefore proves
that its unique negative inverse is $-QX$.
Expanding $3kQX$ using (E96.16) is exactly (E96.7),
which proves (E96.8).

For integrality, the residue field of every such
$K$-prime contains $\mu_3$, hence $k\in3\mathbb Z_3$.
Every coefficient in (E96.7) is integral:
$h_*,c_3,c_\zeta$ are integral, the three path columns
become integral after multiplication by three,
and $n=k/3$ is integral. The quotient $Q$ is integral
as just proved. This proves the asserted numerator
lattice without replacing $M/\mathbb Z_3w$ by $M_0$.
In particular the first row of $Q$ retains the
coefficient $-9$ at the pivot column.

\emph{Global completeness and convergence.}
All the Kummer cochains in (E96.6) are continuous
with values in the indicated integral lattices.
The inertia vectors tend to zero modulo every
power of three, and the other coefficients are
bounded in $\mathbb Z_3$. Formula (E96.7) is therefore
an integral restricted sequence. Equation (E96.8)
identifies it with the very corrections appearing
in E37.11. Division by $3(q-1)$ gives precisely the
two exponents in (E96.9), including the variation
of $v_3(q-1)$ with the prime.

E.37 proves both directions of the global extension
criterion through its continuous localization complex;
restriction from $\mathbb Q$ to $K$ is injective on
the obstruction because restriction followed by
corestriction is multiplication by two. Hence the
set in (E96.9) is the entire global image on the
specified slice. Changing an integral path representative
adds the differential of an integral restricted
sequence, as in E.37. The same applies to the
representative of the abelian cocycle. Thus the
description is on cohomology classes with the displayed
coordinates, rather than on an unrecorded choice of
local representatives. For a valuation vector with
finite support, $v=0$ makes the entire column vanish
at every other prime, proving the finite-support
assertion. The degree-zero stabilizer vanishes by
E.36, so the already evaluated central cohomology
gives exactly (E96.13).

\emph{The complete first inertia obstruction.}
Let $x=X_p$ be the pivot coefficient of (E96.17).
The $f_1\wedge f_2$ coefficient of $X$ is zero,
because every term has a toric factor. The same
coefficient of $w$ is nine. Thus integrality of
$QX$ gives $9x\in\mathbb Z_3$. Put
$\theta=9x\bmod3$.

Write $(u_a,v_a^{\rm pt})$ for the actual coordinates
of the cusp $a$ in E37.2. The compact part of $h_*$
is $(1,0)$, and the retained path columns satisfy
\[
 3\pi_E(b_a)\equiv(u_a,v_a^{\rm pt})\pmod3.
 \tag{E96.19}
\]
This follows from the proper Kummer class of the
actual point $a$; it is also checked for all eight
literal columns. The coefficients of $w$ at
$f_1\wedge e_a$ and $f_2\wedge e_a$ are
$-3u_a$ and $-3v_a^{\rm pt}$, respectively.
Multiplying the two corresponding coordinates of
$QX$ by three and reducing modulo three gives
\[
 (\epsilon-u_a)\overline v_a+\theta u_a=0,
 \qquad
 v_a^{\rm pt}(\overline v_a-\theta)=0.
 \tag{E96.20}
\]
These equations retain the pivot contribution.

For $\epsilon=0$, every cusp has a nonzero point
coordinate, so (E96.20) gives $\overline v_a=\theta$
for all eight cusps. For $\epsilon=1,2$, use the
cusp $B_1$ with point coordinates $(1,1)$: its two
equations give $\epsilon\theta=0$, hence $\theta=0$.
All six cusps with $v_a^{\rm pt}\ne0$ then have
$\overline v_a=0$. The remaining two coordinates
are $A_0=(1,0)$ and $B_0=(2,0)$; their first equations
leave exactly the coordinate asserted in (E96.12).
This proves the necessary image inclusion.

For the reverse image inclusion, take the actual
rational tangent at $D_\epsilon$. Its path Kummer
class has compact coordinate $\epsilon t_2$ and is
unramified away from three. The local geometric
inertia at that tangent is $\mathbb Z_3(1)$, with
the actual cyclotomic Galois action. Twisting its
section by the continuous cocycle $\kappa(a)$ for
$a\in\mathbb Q^\times$ gives a genuine section of
the full geometric group, before taking $N_2$.
After transport to the base tangent its abelian
class is $k_{D_\epsilon}+\kappa(a)e_{D_\epsilon}$.
Its good-prime inertia vector is
$v_p(a)e_{D_\epsilon}$.

Every finitely supported family in
$\bigoplus_{p\ne3}\mathbb F_3$ is obtained from the
finite rational product $a=\prod p^{a_p}$, with
$a_p\in\{0,1,2\}$. Since $K/\mathbb Q$ is unramified
away from three, the valuation at each prime above
$p$ is the same. The cusp directions in (E96.10)
are fixed by the actual cusp permutation, so this
construction has precisely the required conjugate-prime
equivariance. It realizes every family in (E96.11).
For $\epsilon=0$, the sign in $e_O=-\sum e_d$ is
absorbed by replacing $a$ with $a^{-1}$.
This proves equality of the images.

The class $\kappa(7)e_{A_0}$ has compact residue
zero and has inertia vector $e_{A_0}$ at a prime
over seven, which fails the first line of (E96.12).
The certificate also evaluates its full correction
at the actual marked norm-seven Frobenius, and
finds thirty-six nonintegral coordinates. For the
two rational cusp twists, the same local defect
is $k_d(F)\wedge e_d-k_d(F)\wedge e_d=0$;
their forty-four corrections are all zero.

The program retains the exact quotient and its
section, all fixed columns, and all forty-four
expanded numerator polynomials. Their $569$ terms
have denominators prime to three. Seventy-two
complete matrix comparisons agree with the original
$44$-dimensional inverse. These additional matrix
inputs are explicitly labelled algebraic samples;
the separate norm-seven witnesses use the actual
marked arithmetic values. The finite calculation
checks all $59{,}049$ pairs of an inertia vector
and pivot residue across the three strata, giving
exactly three solutions on each stratum. The
universal identities and the global proof above
establish the result for the full continuous groups.
\end{proof}
